% docs/physics/latex/03_nuclear_elastic.tex - SIMC++ physics note 03: A(e,e')A nuclear elastic scattering.
% Source of docs/physics/pdf/SIMCpp_Physics_03_nuclear_elastic.pdf (copied to src/physics/nuclear/ by the build).
% Every displayed formula is a feq/feqa environment citing the Fortran line it transcribes. The event.f lines of this
% reaction are those of upstream commit 20d46fc (written event.f@20d46fc), because upstream 59ccdf6 lost them (U-A1).
\documentclass{simcnote}
\simccommon
\simcnotemeta{03}{A(e,e$'$)A elastic}{A(e,e$'$)A nuclear elastic scattering}%
  {Reaction \texttt{nuclear\_elastic}, physics model \texttt{nuc\_elastic}: elastic scattering from the whole nucleus,
   the charge form factor, radiation and outputs}

\begin{document}
\simcnotemaketitle
\simcabstract{Elastic scattering of the electron from the nucleus as a whole, with the recoiling nucleus as the
``hadron''. The event is generated like H(e,e$'$p) (only the electron direction is drawn), with the nucleus mass in the
elastic kinematics; the cross section is the Mott cross section times $Z^2F^2(Q^2)$ and the recoil factor, with a
parametrised charge form factor and no magnetic term. The reaction was added to SIMC in upstream commit 4df7351 and
lost from event.f in e4dd175; SIMC++ runs the code of commit 20d46fc in both modes.}
\simcnotetoc

% =====================================================================================================================
\section{The process and what SIMC simulates}\label{sec:process}
The reaction is $e + A \to e' + A$ with the nucleus left in its ground state. It is selected with the input flag
\code{doing\_nuc\_elast} = 1 \fsrc{dbase.f:1007, 1155}, which sets the hadron mass to the target mass and turns on
\code{doing\_eep}, without any of \code{doing\_hyd\_elast}, \code{doing\_deuterium}, \code{doing\_heavy}
\fsrc{dbase.f:216-219}:
\begin{feq}{dbase.f:217, 228, 272}
M_h = M_A = \text{\code{targ\%mass\_amu}}\cdot u,\qquad M_h^2 = M_h\cdot M_h.
\label{eq:mh}
\end{feq}
Because \code{doing\_eep} is true the struck mass is still $M_\mathrm{struck}=M_p$ and $M_\mathrm{rec,struck}=0$
\fsrc{dbase.f:297-300}; \code{targ\%Mrec} is the input \code{mrec\_amu}, replaced by $M_A-M_p$ if it is more than 100 MeV
away \fsrc{dbase.f:286-292} (the shipped example \code{ee\_calcium48.inp} gives \code{mrec\_amu} = \code{mass\_amu} and
is replaced, \finding{U-D1}). These two masses are used only in the limits (\secr{sec:gen}, \secr{sec:rad}) and in the
reconstructed $W$; the kinematics use $M_A$. The \code{.hist} title is ``A(e,e$'$)A'' \fsrc{simc.f:682-683}.

\begin{simcwarning}
\textbf{Upstream status.} In upstream SIMC 59ccdf6 the \code{doing\_nuc\_elast} branches of event.f, NtupleInit.f and
results\_write.f are missing (removed by e4dd175, which rebuilt the files from a copy older than 4df7351): the electron
energy is never computed, every try fails in \code{generate\_rad} and the run does not end (\finding{U-A1}). SIMC++
implements the reaction as in commit 20d46fc, in the default and in the \code{-{}-fortran-compat} mode, with the (e,e$'$p)
tree layout. The event.f lines cited in this note are those of 20d46fc (\fsrc{event.f@20d46fc:line}); nuclear\_elastic.f,
radc.f, init.f and dbase.f are cited at 59ccdf6 (unchanged for this reaction).
\end{simcwarning}

What is simulated: the elastic kinematics with the nuclear mass, the charge form factor of \secr{sec:xsec}, the beam,
target, radiation and spectrometers of note 00, with the recoil nucleus transported through the hadron arm. What is not:
magnetic scattering ($W_1=0$, correct only for spin-0 nuclei), any nuclear-structure model beyond the two-parameter form
factor, excited states, the quasi-elastic background, Coulomb distortion of the cross section beyond the focusing factor
of note 00 (no DWBA phase shifts; the form-factor minima are not filled), and the Coulomb or energy-loss effects on a slow
recoil nucleus beyond the generic energy-loss model (which is not tuned for heavy ions). SIMC++ refuses
\code{doing\_nuc\_elast} together with another reaction flag or \code{doing\_phsp} (\finding{U-DEV1}).

% =====================================================================================================================
\section{Kinematics, generated variables and generation volume}\label{sec:gen}
As for H(e,e$'$p), only the electron slopes are generated, flat in the electron edges
(\fsrc{event.f@20d46fc:276-277}); the
hadron slopes and the energies are not (the conditions of \fsrc{event.f@20d46fc:280, 289, 300} do not include
\code{doing\_nuc\_elast}). The vertex limits and the energy ranges are those of hydrogen elastic:
\begin{feqa}{init.f:321-325, 374-379, 438-440, 474-476, 500-503}
E_m^\mathrm{V} &= P_m^\mathrm{V} = M_\mathrm{rec}^\mathrm{V} = T_\mathrm{rec}^\mathrm{V} = 0,\qquad \mathrm{sumEgen}=[0,0],\nonumber\\
\mathrm{gen}^E_e &= [\mathrm{edge}^{E,\mathrm{min}}_e,\ \mathrm{edge}^{E,\mathrm{max}}_e+E^\mathrm{max}_{\gamma2}],\qquad
\mathrm{gen}^E_p = [\mathrm{edge}^{E,\mathrm{min}}_p,\ \mathrm{edge}^{E,\mathrm{max}}_p+E^\mathrm{max}_{\gamma3}],
\label{eq:ranges}
\end{feqa}
where $\mathrm{edge}^E_p$ is computed with $M_h = M_A$ (\ceqr{eq:edges}). The generation volume is the electron solid
angle,
\begin{feq}{simc.f:383, 390 (doing\_nuc\_elast is not in the list of l.394)}
V_\mathrm{gen} = \Delta\Omega_e\quad[\mathrm{sr}].
\label{eq:genvol}
\end{feq}

\paragraph{Photon limit.} The energy-conservation limit of the radiated photons uses the struck mass $M_p$ but the hadron
edge with the nuclear mass:
\begin{feq}{init.f:411-413, 421-425}
E_\gamma^\mathrm{tot,max} = E_\mathrm{beam}^\mathrm{max} + M_p - \mathrm{edge}^{E,\mathrm{min}}_e - \mathrm{edge}^{E,\mathrm{min}}_p,
\qquad \mathrm{edge}^{E,\mathrm{min}}_p\approx M_A,
\label{eq:egammatot}
\end{feq}
which is negative by about $M_A-M_p$ for every nucleus heavier than the proton. See \secr{sec:rad}: without
\code{Egamma\_gen\_max} no radiative event can be generated.

% =====================================================================================================================
\section{Completing the event}\label{sec:complete}
\code{complete\_ev} treats \code{doing\_nuc\_elast} as \code{doing\_hyd\_elast} with $M_h=M_A$
\fsrc{event.f@20d46fc:507, 517, 546, 937}:
\begin{feqa}{event.f@20d46fc:517-523, 546-562, 937-938}
E' &= \frac{E_\mathrm{in}\,M_A}{M_A+E_\mathrm{in}(1-\cos\theta_e)}\quad(\text{fail if }E'>E_\mathrm{in}),\qquad
\hat p = \hat q,\qquad P_A = |\vec q|,\qquad E_A = \sqrt{P_A^2+M_A^2},\nonumber\\
E_m &= P_m = M_\mathrm{rec} = T_\mathrm{rec} = 0,
\label{eq:kin}
\end{feqa}
with $\nu$, $Q^2$, $\vec q$ of \ceqr{eq:nuq2}, the recoil slopes from \code{spectrometer\_angles} and
$\delta_A = 100(P_A-P_p^0)/P_p^0$. For a heavy nucleus $E'\approx E_\mathrm{in}$ and the recoil has a small momentum
$|\vec q|\approx2E_\mathrm{in}\sin(\theta_e/2)$ at an angle near $90^\circ$ (in the shipped example, $\theta_e = 4.51^\circ$,
$E_\mathrm{beam}=2118$ MeV: $|\vec q|\approx167$ MeV/$c$ at $87.6^\circ$, the input \code{spec\%p\%P} and
\code{spec\%p\%theta}).

The reconstruction uses the target mass in the missing energy \fsrc{event.f@20d46fc:1344-1346}:
\begin{feq}{event.f@20d46fc:1344-1346}
T^\mathrm{rec}_\mathrm{rec} = 0,\qquad E_m^\mathrm{rec} = \nu^\mathrm{rec}+M_A-E_A^\mathrm{rec},
\label{eq:recem}
\end{feq}
while $W^\mathrm{rec}$ uses $M_p$ (\ceqr{eq:w}) like every reaction.

% =====================================================================================================================
\section{Jacobians}\label{sec:jac}
\begin{feq}{event.f@20d46fc:498, 1013-1014}
J = \bigl(1+x_e'^2+y_e'^2\bigr)^{-3/2},
\label{eq:jac}
\end{feq}
the electron solid-angle factor only (the hadron factor is applied for the reactions listed in
\fsrc{event.f@20d46fc:1019-1023}, not this one).

% =====================================================================================================================
\section{Cross-section model}\label{sec:xsec}
$S=1$ \fsrc{event.f@20d46fc:1399-1401} and the cross section is \code{sig\_nuc\_elastic}, evaluated at the vertex (and at
the reconstructed event for \code{sigcc\_recon}) \fsrc{event.f@20d46fc:1457-1459}. The code names its source as the
ASTRUC routine of the Mo--Tsai program \fsrc{nuclear\_elastic.f:4}:
\begin{feqa}{nuclear\_elastic.f:16-22, 40-41}
\sigma &= \bigl(W_2+2W_1\tan^2(\theta_e/2)\bigr)\,\sigma_\mathrm{Mott}(E',\theta_e,Q^2)\,\frac{M_A}{M_A+E_\mathrm{in}(1-\cos\theta_e)}
\quad[\mu\mathrm{b/sr}],\nonumber\\
W_1 &= 0,\qquad W_2 = \bigl(Z\,F(Q^2)\bigr)^2,
\label{eq:signuc}
\end{feqa}
with $\sigma_\mathrm{Mott}$ the point-charge formula of note 01 (\code{sigMott}, physics\_proton.f:186-187). The recoil
factor equals $E'/E_\mathrm{in}$ for the elastic $E'$ of \eqr{eq:kin}, so \eqr{eq:signuc} is the usual
$\sigma = Z^2F^2\,\sigma_\mathrm{Mott}E'/E_\mathrm{in}$ of a spin-0 nucleus. The form factor is
\begin{feqa}{nuclear\_elastic.f:53-64}
R &= 1.07\,A^{1/3}\ \mathrm{fm}\quad(\text{except }A=205: 5.470,\ 56: 3.729,\ 28: 3.085,\ 27: 3.035\ \mathrm{fm}),\nonumber\\
F(Q^2) &= \frac{\exp\bigl(-Q^2a^2/6\bigr)}{1+Q^2R^2/6}\quad\bigl(Q^2a^2/6<80;\ 0\text{ otherwise}\bigr),\nonumber\\
a &= 2.4\ \mathrm{fm},\qquad Q^2\ \text{in fm}^{-2} = Q^2[\mathrm{MeV^2}]/(\hbar c)^2,
\label{eq:fgauss}
\end{feqa}
$A$ being compared with \code{nint(A)}. The radii of the four special cases are attributed to de Vries et al.\
\cite{devries} \fsrc{nuclear\_elastic.f:55-61}. The function is called \code{Fgauss}, but only the $A$-independent factor
$\exp(-Q^2a^2/6)$ is a Gaussian; the radius $R$ enters through a monopole factor.

\begin{simcfinding}
\textbf{Radius of the form factor.} At small $Q^2$, $F\simeq1-Q^2(R^2+a^2)/6$, i.e.\ the charge distribution of
\eqr{eq:fgauss} has $\langle r^2\rangle = R^2 + a^2 = R^2 + 5.76$ fm$^2$. If the special values of $R$ are meant as the
r.m.s.\ charge radii of the de Vries compilation that the code cites (for example 3.729 fm for $^{56}$Fe), the form
factor corresponds to an r.m.s.\ radius $\sqrt{R^2+a^2}$ (4.44 fm for $^{56}$Fe), and for the default $R=1.07A^{1/3}$ to
$\sqrt{1.145A^{2/3}+5.76}$ fm (4.57 fm for $A=48$). Whether $R$ is meant as an r.m.s.\ radius, and what the fixed
$a=2.4$ fm stands for, is not documented in the code (nuclear\_elastic.f:53-64). Kept in both modes; reported for the
director, not part of FORTRAN\_DIFFERENCES.md.
\end{simcfinding}

The weight is
\begin{feq}{event.f@20d46fc:1540-1542, 1565-1566}
w = J\cdot w_\mathrm{gen}\cdot\sigma\cdot\bigl(1+V_C/E_\mathrm{beam}\bigr)^2_{[\text{using\_Coulomb}]}\quad[\mu\mathrm{b/sr}],
\label{eq:weight}
\end{feq}
the Coulomb energy $V_C$ being that of note 00, \ceqr{eq:coulomb} (the electron energies are shifted by $V_C$ as for every
reaction). \code{wtcontr} is a number of counts (\ceqr{eq:normfac}). In the example \code{ee\_calcium48} the hadron arm
(HMS, \code{spec\%p\%P} = 166.76 MeV/$c$) is used without its spectrometer model (\code{spect\_mode}) and every event fails
\code{pass\_cuts}, so \code{wtcontr} = 0 (\finding{U-B2}).

% =====================================================================================================================
\section{Radiation, energy loss and decays}\label{sec:rad}
\paragraph{Tail limits.} Tail 1 has its own case, the H(e,e$'$p) formula with the target mass:
\begin{feqa}{radc.f:278-285}
E_b^\mathrm{max,min} &= \frac{M_A\,\mathrm{edge}^{E,\mathrm{max,min}}_e}{M_A-\mathrm{edge}^{E,\mathrm{max,min}}_e(1-\hat u_{e,z})}\quad(E_b^\mathrm{max}=10^{10}\text{ if negative}),\nonumber\\
E_{\gamma1} &\in\bigl[E_\mathrm{in}-E_b^\mathrm{max},\ \min(E_\mathrm{in}-E_b^\mathrm{min},\ \mathrm{edge}^\mathrm{max}_{E_m})\bigr],
\label{eq:tail1}
\end{feqa}
then, as for all reactions, $\min(\cdot,E^\mathrm{max}_{\gamma1})$ and $\pm d_E$ \fsrc{radc.f:306-308}. Tails 2 and 3 follow
\fsrc{radc.f:372-389, 428-446}, with the (e,e$'$p) missing-energy terms void (\code{edge\%Em} = $\pm10^{10}$), and every
upper limit is bounded by $E_\gamma^\mathrm{tot,max}$ (\eqr{eq:egammatot}).

\begin{simcfinding}
\textbf{Negative photon limit.} With \code{doing\_nuc\_elast} the struck mass is $M_p$ (dbase.f:297-300) but the recoil
edge is computed with $M_A$, so $E_\gamma^\mathrm{tot,max}\approx E_\mathrm{beam}+M_p-E'-M_A<0$ (init.f:411-413), and
$E^\mathrm{max}_{\gamma1,2,3}$ are negative. \code{basicrad} then fails every radiated tail (radc.f:61), so with
\code{using\_rad}=1 and \code{Egamma\_gen\_max}=0 no event is ever accepted; the shipped example works only because it sets
\code{Egamma\_gen\_max} = 1000 MeV (\code{hardwired\_rad}, init.f:421). Measured with SIMC++ (\code{-{}-fortran-compat},
mt64 random numbers, seed 4242, \code{ee\_calcium48} with \code{ngen} = $-3000$): \code{Egamma\_max} = $-43296.1$ MeV and 0
successes in 3000 tries with \code{Egamma\_gen\_max} = 0, against 265 successes with 1000 MeV (run folders
\code{simc\_runs/20261123\_171324\_PHYS.1\_nucel\_noEgmax} and \code{\ldots\_171615\_PHYS.1\_nucel\_Egmax1000}). The same
code path is in the 20d46fc Fortran. Kept in both modes; reported for the director.
\end{simcfinding}

\paragraph{Soft and hard factors.} The per-event constants of note 00 (\csecr{sec:radconst}) treat the recoil nucleus as
the third radiating particle: $\lambda_3$ with $E_A$, $P_A$ (\ceqr{eq:lambda}), $c^\mathrm{int}_3$ with $M_p$ in the
normalisation. \code{brem} reconstructs the electron angle and $Q^2$ from \emph{elastic e--p} kinematics
(\ceqr{eq:bremkin}): for a nucleus $E'$ is much closer to $E_\mathrm{in}$ than for a proton at the same angle, so the
angle seen by \code{brem}, $\sin^2(\theta/2)=\frac{M_p}{2E_\mathrm{in}}\bigl(\frac{E_\mathrm{in}}{E'}-1\bigr)$, is smaller
than the true one by about the factor $M_p/M_A$, and $\delta_\mathrm{hard}$ (through $\ln Q^2/m_e^2$) and the e--e
interference term are evaluated at that smaller $Q^2$. This is the general approximation \finding{P5.F1 Q8}; its size is
largest for this reaction. With \code{rad\_proton} the proton terms of \code{brem} use the proton mass, not $M_A$.

\paragraph{Energy loss, multiple scattering, decays.} The recoil nucleus is passed through
\code{trip\_thru\_target} and
\code{target\_musc} with $m=M_A$ (\csecr{sec:eloss}); the energy-loss formula has the ionisation constants of a singly
charged particle (no $z^2$ of the projectile), and the most-probable correction of the reconstruction is applied to it
like to a proton. No decays.

% =====================================================================================================================
\section{Input parameters and output columns}\label{sec:io}
\begin{simctable}{footnotesize}{L{0.32\linewidth}L{0.64\linewidth}}
\caption{Inputs specific to A(e,e$'$)A.}\label{tab:inputs}\\
\toprule\rowcolor{simchead}\simcth{input} & \simcth{role}\\\midrule\endfirsthead
\toprule\rowcolor{simchead}\simcth{input} & \simcth{role}\\\midrule\endhead
\bottomrule\endlastfoot
\code{doing\_nuc\_elast} = 1 & selects the reaction (no other \code{doing\_*} flag)\\
\code{targ\%A}, \code{targ\%Z}, \code{targ\%mass\_amu} & $A$ (radius $R$), $Z$ (form factor, Coulomb), $M_A$\\
\code{spec\%p\%P}, \code{spec\%p\%theta}, \code{hadron\_arm} & the arm that sees the recoil nucleus\\
\code{Egamma\_gen\_max} & must be set ($>0.01$ MeV) with \code{using\_rad} (\secr{sec:rad})\\
\code{using\_Coulomb} & Coulomb energy and focusing factor\\
\end{simctable}
The remaining inputs are those of note 01, \S7.

% docs/physics/latex/frag_eep_columns.tex - the (e,e'p) tree layout, \input by the notes 01, 02 and 03.
% Source: upstream results_write.f:28-313 (59ccdf6) and src/output/NtupleLayout.cpp (UpstreamEepColumns,
% EepNtupleColumns). No displayed equations here (the lint of build_physics_pdf.py checks the notes that input it).
The tree \code{h10} has one entry per successful event \fsrc{simc.f:346} and, in the default (upstream 59ccdf6) layout,
53 columns \fsrc{results\_write.f:42-313}; \code{-{}-tree-layout legacy} gives the 46 columns of SIMC 294f662 (without
\code{xbj}, the vertex block \code{qi}\ldots\code{Pmi}, and with \code{thetapq}, \code{phipq} instead).
Columns 1--12 belong to the arm on the beam right (prefix \code{hs}), 13--24 to the arm on the beam left (\code{ss}),
whatever particle it detects \fsrc{results\_write.f:42-129}.

\begin{simctable}{footnotesize}{L{0.07\linewidth}L{0.22\linewidth}L{0.5\linewidth}L{0.15\linewidth}}
\caption{The (e,e$'$p) columns of the tree \code{h10} (default layout).}\label{tab:eepcolumns}\\
\toprule\rowcolor{simchead}\simcth{\#} & \simcth{name} & \simcth{content} & \simcth{unit}\\\midrule\endfirsthead
\toprule\rowcolor{simchead}\simcth{\#} & \simcth{name} & \simcth{content} & \simcth{unit}\\\midrule\endhead
\bottomrule\endlastfoot
1--4 & \code{hsdelta}, \code{hsyptar}, \code{hsxptar}, \code{hsytar} & reconstructed $\delta$, slopes, $y_\mathrm{tar}$ of the right arm & \%, rad, cm\\
5--8 & \code{hsxfp}, \code{hsxpfp}, \code{hsyfp}, \code{hsypfp} & focal-plane track & cm, rad\\
9--11 & \code{hsdeltai}, \code{hsyptari}, \code{hsxptari} & \code{orig} values (true, after radiation) & \%, rad\\
12 & \code{hsytari} & $z_t\sin\theta_0$ & cm\\
13--24 & \code{ss}\ldots & the same for the left arm (24: $-z_t\sin\theta_0$) & \\
25--28 & \code{q}, \code{nu}, \code{Q2}, \code{W} & reconstructed $|\vec q|$, $\nu$, $Q^2$, $W$ ($W$ with $M_p$) & GeV, GeV$^2$\\
29--32 & \code{xbj}, \code{epsilon}, \code{Em}, \code{Pm} & reconstructed $x$, $\varepsilon$, $E_m$, $P_m$ & --, GeV\\
33--35 & \code{qi}, \code{nui}, \code{Q2i} & vertex $|\vec q|$, $\nu$, $Q^2$ & GeV, GeV$^2$\\
36, 38 & \code{Wi}, \code{epsiloni} & \code{main\%W}, \code{main\%epsilon}: always 0 for (e,e$'$p) and (e,e$'$)A (\finding{U-B1}) & \\
37 & \code{xbji} & vertex $x$ & \\
39--40 & \code{Emi}, \code{Pmi} & vertex $E_\mathrm{miss}=\nu+M_A-E_p$ and $P_\mathrm{miss}$ & GeV\\
41 & \code{corrsing} & $P_e^\mathrm{rec} - M_pE_\mathrm{beam}/(2E_\mathrm{beam}\sin^2(\theta_e^\mathrm{rec}/2)+M_p)$ \fsrc{results\_write.f:33-34} & GeV\\
42--44 & \code{Pmx}, \code{Pmy}, \code{Pmz} & $\vec P_m$ in the frame of \fsrc{results\_write.f:35-39}: $-P_{m,x}$,
  $(P_{m,z}\hat q_y-P_{m,y}\hat q_z)/n$, $-(P_{m,y}\hat q_y+P_{m,z}\hat q_z)/n$, $n=\sqrt{\hat q_y^2+\hat q_z^2}$ & GeV\\
45--47 & \code{PmPar}, \code{PmPer}, \code{PmOop} & reconstructed components of $\vec P_m$ & GeV\\
48 & \code{fry} & $-$\code{main\%target\%rastery} & cm\\
49 & \code{radphot} & $E_{\gamma1}+E_{\gamma2}+E_{\gamma3}$ & GeV\\
50 & \code{sigcc} & $\sigma$ at the vertex (with the Coulomb factor) & see \secr{sec:xsec}\\
51 & \code{Weight} & the event weight $w$ & see \secr{sec:xsec}\\
52, 53 & \code{theta\_e}, \code{theta\_p} & reconstructed polar angles & rad\\
\end{simctable}

For A(e,e$'$)A the hadron columns describe the recoil nucleus, \code{Em} is \eqr{eq:recem}, and \code{corrsing},
\code{Pmx}, \code{Pmy}, \code{Pmz} are 0: results\_write.f computes them only for hydrogen, deuterium and heavy (e,e$'$p)
\fsrc{results\_write.f@20d46fc:54-61}, and the SAVEd locals keep 0 (\finding{U-B2}). \code{sigcc} is $\sigma$ of
\eqr{eq:signuc} in $\mu$b/sr.

% =====================================================================================================================
\section{Fortran \texorpdfstring{$\leftrightarrow$}{<->} C++ map}\label{sec:map}
\begin{simctable}{footnotesize}{L{0.42\linewidth}L{0.55\linewidth}}
\toprule\rowcolor{simchead}\simcth{Fortran} & \simcth{C++ (src/...)}\\\midrule\endfirsthead
\toprule\rowcolor{simchead}\simcth{Fortran} & \simcth{C++ (src/...)}\\\midrule\endhead
\bottomrule\endlastfoot
dbase.f:216-219 (59ccdf6) & \csrc{core/RunSetup.cpp: DeriveReactionFlags} (refuses combinations, U-DEV1)\\
event.f@20d46fc:generate l.276-414 & \csrc{reactions/HydrogenElastic.cpp: HydrogenElastic::generate} (struct \code{NuclearElastic}, reaction \code{nuclear\_elastic})\\
event.f@20d46fc:complete\_ev l.507-563, 937-938 & \csrc{reactions/EepReactionBase.cpp: EepReactionBase::completeEv}\\
event.f@20d46fc:complete\_recon\_ev l.1344-1346 & \csrc{EepReactionBase::completeRecon}\\
event.f@20d46fc:complete\_main l.1399-1401, 1457-1459 & \csrc{HydrogenElastic::spectralWeight}; \csrc{reactions/EventChain.cpp: assembleWeight}\\
nuclear\_elastic.f: sig\_nuc\_elastic, NUC\_FORM\_FACTOR, Fgauss & \csrc{physics/nuclear/NuclearElastic.cpp: SigNucElastic, NucFormFactor, Fgauss}; model \code{nuc\_elastic}\\
radc.f:278-286 (59ccdf6) & \csrc{radiation/Radc.cpp: generateRad} (CASE 2b)\\
results\_write.f@20d46fc:246, NtupleInit.f@20d46fc:293 & \csrc{output/NtupleLayout.cpp: UsesEepLayout, FillEepNtuple, FillUpstreamNtuple}\\
\end{simctable}

% =====================================================================================================================
\section{Fortran findings and the two modes}\label{sec:findings}
\begin{simctable}{footnotesize}{L{0.17\linewidth}L{0.5\linewidth}L{0.29\linewidth}}
\toprule\rowcolor{simchead}\simcth{id} & \simcth{what} & \simcth{default / \code{-{}-fortran-compat}}\\\midrule\endfirsthead
\toprule\rowcolor{simchead}\simcth{id} & \simcth{what} & \simcth{default / \code{-{}-fortran-compat}}\\\midrule\endhead
\bottomrule\endlastfoot
\finding{U-A1} (A) & upstream 59ccdf6 lost the event.f/ntuple branches: endless run with 0 successes & the 20d46fc code in both modes\\
\finding{U-B2} (B) & $W_1=0$, the \code{Fgauss} model, $W$ with $M_p$, \code{corrsing}/\code{Pm*} never computed (0) & kept in both\\
\finding{U-D1} (D) & \code{.hist} line \code{doing\_nuc\_elast}; the example input quirks (\code{Ebeam = 2118.0.}, no seed,
  \code{mrec\_amu} replaced) & followed\\
\finding{U-DEV1} (DEV) & the Fortran lets \code{doing\_nuc\_elast} mix with another reaction & refused at set-up\\
\finding{P5.F1 Q8} (D) & \code{brem} uses e--p elastic kinematics; largest effect here & kept\\
candidate & $E_\gamma^\mathrm{tot,max}<0$ without \code{Egamma\_gen\_max} (\secr{sec:rad}) & kept; reported\\
candidate & $\langle r^2\rangle = R^2+(2.4\ \mathrm{fm})^2$ in \code{Fgauss} (\secr{sec:xsec}) & kept; reported\\
\end{simctable}
The 5000-event \code{ee\_calcium48} run of SIMC++ (RANLUX, \code{-{}-fortran-compat}) is compared with the 20d46fc Fortran
in the tests \code{l1\_simcrun\_ranlux\_up\_ee\_calcium48\_*} (FORTRAN\_DIFFERENCES.md, U-A1).

\begin{thebibliography}{9}
\bibitem{devries} H.~de Vries, C.~W.~de Jager, C.~de Vries, ``Nuclear charge-density-distribution parameters from elastic
  electron scattering'', Atomic Data and Nuclear Data Tables \textbf{36}, 495 (1987) (cited at
  \fsrc{nuclear\_elastic.f:55-57}).
\bibitem{mo} L.~W.~Mo and Y.~S.~Tsai, Rev.\ Mod.\ Phys.\ \textbf{41}, 205 (1969) (``ASTRUC of Mo-Tsai program'',
  \fsrc{nuclear\_elastic.f:4}).
\end{thebibliography}
\end{document}
